\chapter{Systematyczny przegląd literatury}
\label{chapter:slr}

\section{Cel i zakres przeglądu}

Częściowy systematyczny przegląd literatury przeprowadzono w celu uporządkowania aktualnego stanu wiedzy dotyczącego porównania architektury monolitycznej i mikroserwisowej w środowiskach chmurowych. Analizie poddano publikacje dotyczące wpływu stylu architektonicznego na wydajność, skalowalność oraz koszty utrzymania systemów, z uwzględnieniem czynników kontekstowych, takich jak konteneryzacja, orkiestracja, obserwowalność i model komunikacji między usługami.

Przegląd literatury pełni w opisywanej pracy podwójną funkcję. Po pierwsze, pozwala określić, jakie podejścia badawcze dominują w dotychczasowych opracowaniach i jakie metryki są najczęściej wykorzystywane do oceny obu architektur. Po drugie, umożliwia identyfikację luk badawczych uzasadniających przeprowadzenie własnego eksperymentu porównawczego.


\section{Pytania badawcze}

Na potrzeby przeglądu sformułowano główne pytanie badawcze:

\begin{quote}
    \textbf{Q:} W jaki sposób wybór pomiędzy architekturą monolityczną a architekturą opartą na mikroserwisach wpływa na wydajność, skalowalność oraz koszty utrzymania aplikacji wdrożonej w środowisku chmurowym?
\end{quote}

W celu uporządkowania analizy główne pytanie badawcze rozbito na cztery pytania szczegółowe:

\begin{itemize}
    \item \textbf{Q1:} Jak architektura monolityczna i mikroserwisowa wypadają w porównaniu pod względem wydajności, mierzonej za pomocą takich wskaźników jak czas odpowiedzi, latencja, przepustowość oraz dostępność?
    \item \textbf{Q2:} Jakie różnice występują pomiędzy obiema architekturami w zakresie skalowalności pionowej i poziomej oraz zachowania systemu przy rosnącym obciążeniu?
    \item \textbf{Q3:} Jakie koszty utrzymania i koszty operacyjne są raportowane dla architektury monolitycznej i mikroserwisowej, w szczególności w odniesieniu do infrastruktury, utrzymania, obserwowalności oraz procesów DevOps?
    \item \textbf{Q4:} W jakich warunkach technologicznych i środowiskowych, takich jak rodzaj aplikacji, użyty stos technologiczny, środowisko chmurowe czy wykorzystanie kontenerów i orkiestracji, wyniki porównań ulegają zmianie?
\end{itemize}


\section{Metodyka przeglądu}

Przegląd przeprowadzono zgodnie z podejściem systematycznym, obejmującym zdefiniowanie pytania badawczego, opracowanie strategii wyszukiwania, wybór źródeł, zastosowanie kryteriów włączenia i wykluczenia, a następnie ekstrakcję oraz syntezę danych. Jako ramę konceptualną przyjęto model \textbf{PICO}, który dobrze odpowiada charakterowi porównawczemu niniejszego badania:

\begin{itemize}
    \item \textbf{P (Population):} systemy i aplikacje wdrażane w środowiskach chmurowych,
    \item \textbf{I (Intervention):} architektura mikroserwisowa,
    \item \textbf{C (Comparison):} architektura monolityczna,
    \item \textbf{O (Outcome):} wydajność, skalowalność, koszty utrzymania i utrzymywalność.
\end{itemize}


\section{Strategia wyszukiwania}

Strategia wyszukiwania została oparta na czterech grupach pojęć logicznie związanych z tematem pracy: (1) style architektoniczne, (2) środowiska chmurowe i uruchomieniowe, (3) metryki techniczne oraz (4) koszty i utrzymanie. Wewnątrz każdej grupy stosowano operator logiczny \texttt{OR}, natomiast pomiędzy grupami operator \texttt{AND}. Dzięki temu zapytanie obejmowało publikacje rzeczywiście odnoszące się do porównania monolitu i mikroserwisów w kontekście chmurowym.

W praktyce wykorzystano następujące grupy słów kluczowych:

\begin{itemize}
    \item \textbf{Architektura:} \texttt{microservice*}, \texttt{"micro-service*"}, \texttt{"microservices architecture"}, \texttt{"microservice architecture"}, \texttt{"service-oriented architecture"}, \texttt{monolith*}, \texttt{"monolithic architecture"}, \texttt{"monolithic application"},
    \item \textbf{Chmura i runtime:} \texttt{cloud}, \texttt{"cloud computing"}, \texttt{"public cloud"}, \texttt{"cloud-native"}, \texttt{kubernetes}, \texttt{container*}, \texttt{docker}, \texttt{orchestration},
    \item \textbf{Metryki techniczne:} \texttt{performance}, \texttt{benchmark*}, \texttt{throughput}, \texttt{latency}, \texttt{"response time"}, \texttt{scalab*}, \texttt{autoscal*}, \texttt{elasticity},
    \item \textbf{Koszty i utrzymanie:} \texttt{cost}, \texttt{"operational cost"}, \texttt{"infrastructure cost"}, \texttt{TCO}, \texttt{"total cost of ownership"}, \texttt{maintenance}, \texttt{maintainability}, \texttt{reliability}, \texttt{availability}, \texttt{observability}, \texttt{DevOps}.
\end{itemize}

Główne zapytanie wyszukiwawcze miało następującą postać:
\begin{listing}[ht]
    \begin{minted}[fontsize=\small, breaklines]{text}
(
  (microservice* OR "micro-service*" OR "microservices architecture" OR "microservice architecture" OR "service-oriented architecture")
  AND
  (monolith* OR "monolithic architecture" OR "monolithic application" OR monolith)
  AND
  (cloud OR "cloud computing" OR "public cloud" OR "cloud-native" OR kubernetes OR container* OR docker OR orchestration)
  AND
  (performance OR benchmark* OR throughput OR latency OR "response time" OR scalab* OR autoscal* OR elasticity)
  AND
  (cost OR "operational cost" OR "infrastructure cost" OR TCO OR "total cost of ownership" OR maintenance OR maintainability OR reliability OR availability OR observability OR DevOps)
)
AND publication year > 2016
AND document type IN (article, conference paper)
AND language = English
    \end{minted}
    \caption{Zapytanie do wyszukiwania artykułów}
    \label{lst:slr-search-sequence}
\end{listing}

\section{Źródła danych i zakres przeszukiwania}

Wyszukiwanie literatury przeprowadzono w bazach uznawanych za podstawowe źródła publikacji z obszaru informatyki i inżynierii oprogramowania. W szczególności uwzględniono Scopus~\cite{scopus}, IEEE Xplore~\cite{ieeexplore}, Web of Science~\cite{webofscience}. Dobór źródeł miał zapewnić zarówno szerokie pokrycie literatury, jak i dostęp do publikacji szczególnie informatycznych oraz konferencyjnych.

Uwzględniono publikacje opublikowane w języku angielskim po roku 2016. Przyjęcie takiego zakresu czasowego było uzasadnione faktem, że po 2016 roku nastąpił wyraźny wzrost liczby badań dotyczących systemów mikroserwisowych, konteneryzacji i cloud-native. Wyszukiwanie obejmowało artykuły naukowe oraz materiały konferencyjne, natomiast pominięto publikacje o charakterze czysto marketingowym, notki redakcyjne oraz opracowania bez wartości empirycznej.

Łącznie w przeszukiwanych bazach zidentyfikowano 469 rekordów. Rozkład liczby rekordów pozyskanych z poszczególnych baz przedstawiono w tabeli~\ref{tab:slr-databases}. Należy zaznaczyć, że bazy wykorzystane do wyszukiwania literatury nie są tożsame ze źródłami publikacyjnymi widocznymi w końcowym zbiorze. Bazy takie jak Scopus czy Web of Science agregują publikacje pochodzące od różnych wydawców i platform publikacyjnych.

\begin{table}[htbp]
    \centering
    \caption{Liczba rekordów pozyskanych z przeszukiwanych baz literaturowych}
    \label{tab:slr-databases}
    \begin{tabular}{|l|r|}
        \hline
        \textbf{Baza literaturowa} & \textbf{Liczba rekordów} \\
        \hline
        Scopus & 199 \\
        Web of Science & 135 \\
        IEEE Xplore & 135 \\
        \hline
        \textbf{Łącznie} & \textbf{469} \\
        \hline
    \end{tabular}
\end{table}

\section{Kryteria włączenia i wyłączenia}

W celu zapewnienia spójności procesu selekcji zdefiniowano kryteria włączenia i wyłączenia.

\subsection{Kryteria włączenia}

Do dalszej analizy kwalifikowano publikacje, które spełniały co najmniej jedno z poniższych kryteriów:
\begin{itemize}
    \item porównywały architekturę mikroserwisową i monolityczną bezpośrednio lub zawierały dane umożliwiające takie porównanie,
    \item dotyczyły systemów uruchamianych w środowiskach chmurowych lub kontenerowych,
    \item raportowały wyniki odnoszące się do wydajności, skalowalności, utrzymywalności, kosztów lub operacyjności,
    \item miały charakter eksperymentalny, benchmarkingowy lub zawierały systematyczną analizę literatury istotną dla przyjętego problemu badawczego.
\end{itemize}

\subsection{Kryteria wyłączenia}

Z analizy wyłączano publikacje, które:
\begin{itemize}
    \item nie dotyczyły bezpośrednio architektury monolitycznej i mikroserwisowej,
    \item nie odnosiły się do środowisk chmurowych, kontenerowych ani nowoczesnych środowisk wdrożeniowych,
    \item były duplikatami rekordów pozyskanych z innych baz.
\end{itemize}


\section{Proces selekcji publikacji}

Proces selekcji miał charakter dwuetapowy. W pierwszym etapie przeprowadzono screening tytułów i abstraktów, odrzucając publikacje jednoznacznie niezwiązane z problemem badawczym. W drugim etapie analizowano pełne teksty prac zakwalifikowanych wstępnie, sprawdzając ich zgodność z pytaniami badawczymi oraz przyjętymi kryteriami selekcji.

Do organizacji procesu screeningu i usuwania duplikatów wykorzystano narzędzie Rayyan, które umożliwia sprawne zarządzanie rekordami bibliograficznymi, ich tagowanie oraz wspomaga deduplikację. Dodatkowo publikacjom przypisywano znaczniki tematyczne odpowiadające głównym osiom analizy, takim jak \textit{performance}, \textit{scalability}, \textit{cost}, \textit{Kubernetes}, \textit{benchmark} czy \textit{empirical study}.

Kolejne etapy zawężania zbioru publikacji przedstawiono w tabeli~\ref{tab:slr-selection-steps}. Proces obejmował usunięcie duplikatów, selekcję na podstawie tytułu i abstraktu, analizę pełnych tekstów oraz ocenę według kryteriów jakościowych.

\begin{table}[htbp]
    \centering
    \caption{Etapy selekcji publikacji w systematycznym przeglądzie literatury}
    \label{tab:slr-selection-steps}
    \begin{tabular}{|l|r|r|r|}
        \hline
        \textbf{Etap} & \textbf{Liczba przed etapem} & \textbf{Liczba po etapie} & \textbf{Liczba odrzuconych} \\
        \hline
        Usunięcie duplikatów & 469 & 390 & 79 \\
        Selekcja po tytule i abstrakcie & 390 & 39 & 351 \\
        Selekcja po pełnej treści & 39 & 24 & 15 \\
        Selekcja po kryteriach jakościowych & 24 & 16 & 8 \\
        \hline
    \end{tabular}
\end{table}

Po zakończeniu procesu deduplikacji i wieloetapowej selekcji końcowy zbiór obejmował 16 publikacji. Jak pokazano w tabeli~\ref{tab:slr-selection-steps}, największa liczba rekordów została odrzucona na etapie selekcji po tytule i abstrakcie, co wskazuje, że szeroka strategia wyszukiwania skutecznie zapewniła wysoką czułość wyszukiwania, kosztem konieczności późniejszego odfiltrowania publikacji mniej istotnych dla problemu badawczego. Rozkład publikacji według roku wydania wskazuje, że największa liczba prac została opublikowana w 2023 roku, natomiast wyraźny wzrost zainteresowania tematyką jest widoczny od 2020 roku (rys.~\ref{fig:slr-years}).

Znacząca część końcowego zbioru obejmowała publikacje pochodzące ze źródeł publikacyjnych związanych z IEEE Xplore~\cite{ieeexplore}, które stanowiły 56{,}25\% wszystkich zakwalifikowanych prac. Drugim najważniejszym źródłem był SpringerLink~\cite{springerlink} z udziałem 18{,}75\%, natomiast ACM Digital Library~\cite{acm_dl}, ScienceDirect~\cite{sciencedirect}, MDPI~\cite{mdpi} oraz CEUR-WS~\cite{ceur} odpowiadały po 6{,}25\% zbioru. Taki rozkład wskazuje, że analizowana literatura ma głównie charakter informatyczny i konferencyjny, z wyraźną reprezentacją źródeł technicznych (rys.~\ref{fig:slr-sources}).

\begin{figure}[htbp]
    \centering
    \includegraphics[width=0.82\textwidth]{images/slr/wykres_publikacje_lata}
    \caption{Liczba wybranych publikacji w zależności od roku wydania}
    \label{fig:slr-years}
\end{figure}

\begin{figure}[htbp]
    \centering
    \includegraphics[width=0.82\textwidth]{images/slr/wykres_zrodla_publikacji}
    \caption{Udział poszczególnych źródeł publikacyjnych w końcowym zbiorze}
    \label{fig:slr-sources}
\end{figure}


\section{Ekstrakcja i synteza danych}

Dla każdej zakwalifikowanej publikacji przeprowadzono ekstrakcję danych według jednolitego schematu. Rejestrowano przede wszystkim:
\begin{itemize}
    \item typ badania,
    \item analizowany styl architektoniczny,
    \item kontekst wdrożeniowy,
    \item zastosowane technologie,
    \item metryki oceny,
    \item główne wyniki oraz ograniczenia badania.
\end{itemize}

W celu uporządkowania materiału źródłowego zastosowano syntezę tematyczną. Publikacje podzielono na trzy główne grupy odpowiadające celom niniejszej pracy:
\begin{enumerate}
    \item badania koncentrujące się na wydajności,
    \item badania dotyczące skalowalności i zachowania pod obciążeniem,
    \item badania odnoszące się do kosztów utrzymania, utrzymywalności oraz złożoności operacyjnej.
\end{enumerate}


\section{Analiza badań powiązanych dotyczących porównania architektur}

Analizowane publikacje nie wskazują jednoznacznej przewagi architektury monolitycznej albo mikroserwisowej. Wyniki zależą od badanego scenariusza, charakteru obciążenia, języka programowania, sposobu podziału systemu, środowiska uruchomieniowego oraz przyjętych metryk. Z tego względu szczególne znaczenie mają badania, w których porównywano co najmniej dwa warianty architektoniczne tej samej lub funkcjonalnie zbliżonej aplikacji. Pozwala to ograniczyć liczbę zmiennych zakłócających i lepiej ocenić wpływ samego stylu architektonicznego.

\subsection*{Badania eksperymentalne nad wydajnością i skalowalnością}

\noindent\textbf{Tapia i in. -- \textit{From Monolithic Systems to Microservices: A Comparative Study of Performance}.}
Tapia i in.~\cite{tapia2020monolithic} przeprowadzili porównanie aplikacji monolitycznej oraz jej wariantu mikroserwisowego, koncentrując się na relacji między architekturą systemu a wykorzystaniem zasobów. W badaniu analizowano między innymi zużycie procesora, pamięci, operacje dyskowe oraz ruch sieciowy. Istotne jest jednak to, że wariant monolityczny uruchomiono na maszynie wirtualnej KVM~\cite{kvm}, natomiast wariant mikroserwisowy w kontenerach. Oznacza to, że uzyskane wyniki należy interpretować ostrożnie, ponieważ różnice mogą wynikać zarówno z architektury, jak i ze zmiany środowiska uruchomieniowego. Dla niniejszej pracy publikacja ta jest ważna, ponieważ pokazuje konieczność analizowania nie tylko czasu odpowiedzi lub przepustowości, lecz także narzutu zasobowego związanego z rozproszeniem aplikacji.

\medskip

\noindent\textbf{Blinowski, Ojdowska i Przybyłek -- \textit{Monolithic vs. Microservice Architecture: A Performance and Scalability Evaluation}.}
Blinowski i in.~\cite{blinowski2022monolithic} przedstawili jedno z najbliższych tematycznie badań względem niniejszej pracy. Autorzy porównali architekturę monolityczną i mikroserwisową na przykładzie referencyjnej aplikacji webowej przygotowanej w czterech wariantach: monolit Java, mikroserwisy Java, monolit C\#/.NET oraz mikroserwisy C\#/.NET. Badanie obejmowało scenariusze lokalne oraz chmurowe, a analizowanymi aspektami były między innymi przepustowość, skalowanie pionowe, skalowanie poziome oraz koszt środowiska. Wyniki pracy wskazują, że na pojedynczej maszynie monolit może osiągać lepsze rezultaty niż jego odpowiednik mikroserwisowy, a skuteczność skalowania zależy od charakteru usługi, technologii implementacyjnej oraz konfiguracji infrastruktury. Publikacja ta stanowi bezpośredni punkt odniesienia dla niniejszej pracy, ponieważ również łączy porównanie architektur z analizą wpływu technologii implementacyjnej i sposobu skalowania.

\medskip

\noindent\textbf{Jatkiewicz i Okrój -- \textit{Differences in performance, scalability, and cost of using microservice and monolithic architecture}.}
Jatkiewicz i Okrój~\cite{okroj2023differences} skupili się na porównaniu wydajności, skalowalności i kosztów wykorzystania architektury monolitycznej oraz mikroserwisowej. Autorzy zestawili dwie aplikacje realizujące tę samą logikę biznesową, ale zaimplementowane w odmiennych wariantach architektonicznych. Testy przeprowadzono z użyciem narzędzia JMeter, zarówno w środowisku lokalnym, jak i na platformie Microsoft Azure. Szczególne znaczenie tej pracy wynika z połączenia trzech wymiarów oceny: wydajności, skalowalności oraz kosztu. Z punktu widzenia niniejszej pracy istotny jest wniosek, że mikroserwisy nie powinny być traktowane jako rozwiązanie automatycznie bardziej efektywne. Ich opłacalność zależy od charakteru obciążenia, konfiguracji środowiska oraz sposobu wykorzystania mechanizmów skalowania.

\medskip

\noindent\textbf{Bernsteiner i in. -- \textit{Performance Comparison between a Monolithic and a Microservice Application}.}
Bernsteiner i in.~\cite{bernsteiner2025performance} przeanalizowali migrację wybranych komponentów aplikacji z architektury monolitycznej do mikroserwisowej. Badanie koncentrowało się na porównaniu czasu odpowiedzi oraz użycia procesora. Jest ono istotne, ponieważ dotyczy migracji tej samej aplikacji domenowej, a nie porównania dwóch całkowicie niezależnych systemów. Jednocześnie publikacja pokazuje typowy problem badań migracyjnych: zmiana architektury często pociąga za sobą zmianę frameworków, bazy danych, sposobu komunikacji i środowiska wdrożeniowego. Z tego względu wyniki takiego eksperymentu należy interpretować jako efekt całego procesu migracji, a nie wyłącznie jako efekt przejścia z monolitu do mikroserwisów.

\subsection*{Badania uwzględniające koszt i środowisko chmurowe}

\noindent\textbf{Villamizar i in. -- \textit{Cost comparison of running web applications in the cloud using monolithic, microservice, and AWS Lambda architectures}.}
Villamizar i in.~\cite{villamizar2017cost} porównali koszt uruchamiania aplikacji webowej w chmurze w trzech wariantach: monolitycznym, mikroserwisowym zarządzanym przez użytkownika chmury oraz mikroserwisowym zarządzanym przez dostawcę, w tym AWS Lambda. Publikacja ta jest ważna, ponieważ pokazuje, że ocena architektury nie powinna ograniczać się do samej wydajności, lecz powinna uwzględniać również model wdrożenia i sposób rozliczania zasobów. Autorzy wskazują, że mikroserwisy mogą obniżać koszty infrastruktury względem klasycznej architektury monolitycznej, szczególnie wtedy, gdy wykorzystuje się usługi zarządzane przez dostawcę chmury. Dla niniejszej pracy jest to punkt odniesienia przy interpretowaniu wyników skalowania w kontekście efektywności kosztowej.

\medskip

\noindent\textbf{Mangwani, Mangwani i Motwani -- \textit{Evaluation of a Multitenant SaaS Using Monolithic and Microservice Architectures}.}
Mangwani i in.~\cite{mangwani2023evaluation} przeprowadzili porównanie aplikacji typu multitenant SaaS przygotowanej w wariancie monolitycznym i mikroserwisowym. Aplikację wdrożono jako obrazy Docker w środowisku Google Cloud, a do generowania obciążenia i pomiarów wykorzystano JMeter. Autorzy analizowali między innymi wydajność, skalowalność, równoważenie obciążenia, niezawodność, wykorzystanie zasobów oraz koszt infrastruktury. Publikacja ta ma znaczenie dla niniejszej pracy, ponieważ pokazuje, że potencjalna przewaga mikroserwisów jest silnie związana z możliwością niezależnego skalowania usług. Jednocześnie badanie potwierdza, że mikroserwisy należy oceniać razem ze środowiskiem konteneryzacji i chmury, a nie jako abstrakcyjny wzorzec architektoniczny oderwany od sposobu wdrożenia.

\subsection*{Wnioski dla niniejszej pracy}

Przegląd badań powiązanych wskazuje, że wyniki porównań architektury monolitycznej i mikroserwisowej są silnie zależne od kontekstu eksperymentu. Monolit często osiąga korzystniejsze wyniki w prostszych scenariuszach, zwłaszcza na pojedynczym węźle, ponieważ nie ponosi kosztu komunikacji między usługami. Mikroserwisy mogą natomiast uzyskać przewagę wtedy, gdy aplikacja została poprawnie zdekomponowana, a środowisko umożliwia niezależne skalowanie wybranych komponentów.

Wnioski te uzasadniają przyjęcie w niniejszej pracy porównania obejmującego nie tylko średni czas odpowiedzi, lecz także przepustowość, latencję, sposób skalowania oraz technologię implementacyjną. Z tego względu osobno analizowane są różne style architektoniczne, języki programowania oraz wyniki skalowania wertykalnego i horyzontalnego.

\subsection{Wydajność i narzut komunikacyjny}

W badaniach dotyczących wydajności często wskazuje się, że architektura monolityczna osiąga korzystniejsze wyniki w scenariuszach, w których większość logiki biznesowej wykonywana jest wewnątrz jednego procesu. Wynika to z braku konieczności komunikacji sieciowej pomiędzy usługami, mniejszego narzutu serializacji i deserializacji danych oraz prostszego przepływu żądania przez system. Takie wnioski pojawiają się między innymi w pracach Blinowskiego i in.~\cite{blinowski2022monolithic}, Jatkiewicza i Okroja~\cite{okroj2023differences} oraz Bernsteinera i in.~\cite{bernsteiner2025performance}.

Nie oznacza to jednak, że mikroserwisy są wydajnościowo gorsze w każdym przypadku. W publikacjach analizujących systemy o większej złożoności podkreśla się, że architektura mikroserwisowa może osiągać dobre wyniki wtedy, gdy dekompozycja systemu odpowiada rzeczywistym granicom domenowym, a najbardziej obciążone komponenty mogą być rozwijane i optymalizowane niezależnie. W takim modelu koszt komunikacji między usługami może zostać zrównoważony przez możliwość lepszego rozłożenia obciążenia oraz niezależnego skalowania wybranych części systemu~\cite{tapia2020monolithic,nayim2023ecommerce,adrio2023graphql,kader2025streaming}.

\subsection{Skalowalność i zachowanie pod obciążeniem}

W zakresie skalowalności literatura częściej wskazuje na potencjalną przewagę architektury mikroserwisowej. Jej podstawowym atutem jest możliwość niezależnego skalowania tylko tych komponentów, które faktycznie stają się wąskim gardłem systemu. W przypadku monolitu zwiększanie liczby instancji oznacza replikację całej aplikacji, nawet jeżeli wzrost obciążenia dotyczy jedynie wybranej funkcji. W środowiskach chmurowych i kontenerowych taka różnica może mieć istotne znaczenie dla efektywności wykorzystania zasobów~\cite{blinowski2022monolithic,mangwani2022container,mangwani2023evaluation,okroj2023differences}.

Jednocześnie badania pokazują, że samo przejście na mikroserwisy nie gwarantuje liniowego wzrostu przepustowości. Przy zbyt dużej liczbie instancji lub zbyt intensywnej komunikacji między usługami narzut związany z routingiem, równoważeniem obciążenia, komunikacją sieciową i synchronizacją może ograniczać korzyści ze skalowania poziomego~\cite{blinowski2022monolithic,okroj2023differences,bernsteiner2025performance}. Z tego powodu analiza skalowalności powinna obejmować zarówno skalowanie poziome, jak i pionowe, a także powinna uwzględniać punkt, w którym dalsze zwiększanie zasobów przestaje przynosić proporcjonalny wzrost wydajności.

\subsection{Koszty utrzymania i koszty operacyjne}

Wyniki dotyczące kosztów są mniej jednoznaczne niż wyniki dotyczące wydajności. Część badań analizuje koszt wąsko, jako koszt infrastruktury chmurowej, liczby maszyn wirtualnych, kontenerów lub zużycia zasobów. Inne publikacje traktują koszt szerzej, uwzględniając również złożoność wdrożeń, konieczność utrzymania monitoringu, logowania, automatyzacji CI/CD, konfiguracji środowisk oraz kompetencji zespołu~\cite{villamizar2017cost,papakonstantinou2020technoeconomic,fontanadenardin2021energy}.

Architektura monolityczna może być korzystniejsza kosztowo w małych i średnich systemach, gdzie prostsze wdrożenie, mniejsza liczba komponentów oraz łatwiejsze monitorowanie ograniczają nakład operacyjny. Mikroserwisy mogą natomiast uzasadniać większy koszt organizacyjny i infrastrukturalny wtedy, gdy system osiąga skalę, przy której niezależne wdrażanie, niezależne skalowanie i separacja odpowiedzialności przynoszą realne korzyści~\cite{villamizar2017cost,benavente2022comparative,okroj2023differences,mangwani2023evaluation,berry2024migrating}.

\subsection{Wnioski z analizy badań powiązanych}

Analiza badań powiązanych wskazuje, że dotychczasowe wyniki nie pozwalają na jednoznaczne uznanie jednej architektury za ogólnie lepszą. Monolit częściej uzyskuje korzystne wyniki w prostszych scenariuszach i na pojedynczym węźle, ponieważ nie generuje narzutu komunikacji między usługami. Mikroserwisy mogą natomiast uzyskać przewagę w systemach o większej złożoności, szczególnie wtedy, gdy możliwe jest niezależne skalowanie najbardziej obciążonych komponentów.

Jednocześnie wiele analizowanych badań łączy zmianę architektury ze zmianą środowiska uruchomieniowego, frameworków, mechanizmów komunikacji lub sposobu wdrożenia. Utrudnia to jednoznaczne określenie, czy obserwowane różnice wynikają z samego stylu architektonicznego, czy z pozostałych czynników technologicznych.

\section{Luki badawcze i uzasadnienie własnego podejścia}

Analiza wybranych publikacji wskazuje na kilka istotnych luk badawczych. Po pierwsze, stosunkowo niewiele prac porównuje architekturę monolityczną i mikroserwisową na podstawie funkcjonalnie równoważnych implementacji tej samej aplikacji. W wielu przypadkach porównywane systemy różnią się nie tylko stylem architektonicznym, lecz także zakresem funkcjonalnym, środowiskiem uruchomieniowym lub zastosowanym stosem technologicznym, co utrudnia jednoznaczną interpretację wyników~\cite{tapia2020monolithic,benavente2022comparative,berry2024migrating}.

Po drugie, część badań koncentruje się głównie na pojedynczym wymiarze oceny, najczęściej wydajności albo skalowalności, pomijając szerszy kontekst kosztów utrzymania i złożoności operacyjnej. Tymczasem decyzja o wyborze architektury ma charakter wielokryterialny i powinna uwzględniać zarówno parametry techniczne, jak i konsekwencje infrastrukturalne oraz organizacyjne~\cite{papakonstantinou2020technoeconomic,okroj2023differences,mangwani2023evaluation}.

Po trzecie, w literaturze brakuje jednolitej metodyki eksperymentalnej pozwalającej bezpośrednio porównywać wyniki uzyskane dla różnych architektur, języków programowania i mechanizmów skalowania. Różnice w metrykach, poziomach obciążenia, typach infrastruktury oraz narzędziach testowych ograniczają możliwość prostego zestawiania rezultatów między badaniami.

Zidentyfikowane luki uzasadniają podejście przyjęte w niniejszej pracy. Polega ono na porównaniu funkcjonalnie równoważnych wariantów aplikacji monolitycznej i mikroserwisowej, uruchomionych w porównywalnym środowisku chmurowym i badanych według jednolitej procedury eksperymentalnej. Takie podejście pozwala ograniczyć wpływ czynników zakłócających i lepiej ocenić wpływ stylu architektonicznego na wydajność, skalowalność oraz koszt utrzymania systemu.

\section{Podsumowanie}

Przeprowadzony częściowy systematyczny przegląd literatury potwierdził, że porównanie architektury monolitycznej i mikroserwisowej pozostaje aktualnym problemem badawczym. Analiza wykazała, że monolit może oferować niższy narzut komunikacyjny i większą prostotę operacyjną, natomiast mikroserwisy zapewniają większą elastyczność skalowania oraz lepszą modularność systemu. Ostateczna ocena obu podejść zależy jednak od kontekstu wdrożeniowego, charakterystyki aplikacji, sposobu dekompozycji oraz środowiska uruchomieniowego~\cite{blinowski2022monolithic,okroj2023differences,mangwani2023evaluation,berry2024migrating}.

Wnioski z przeglądu stanowią podstawę do zaprojektowania własnego eksperymentu badawczego, w którym porównane zostaną funkcjonalnie równoważne warianty aplikacji monolitycznej i mikroserwisowej.